% ============================================================
%  CV.tex – Özgeçmiş
% ============================================================
\clearpage
\thispagestyle{empty}

\begin{center}
{\Large \textbf{ÖZGEÇMİŞ}}
\end{center}

\vspace{1cm}

\noindent\textbf{Ad Soyad:} Seyit Ali YORGUN\\[4pt]
\noindent\textbf{Doğum Yeri ve Yılı:} \\[4pt]
\noindent\textbf{E-posta:} \\[4pt]

\vspace{0.8cm}

\noindent\textbf{Öğrenim Durumu:}

\vspace{0.3cm}

\begin{tabular}{@{}lll@{}}
\textbf{Derece} & \textbf{Kurum} & \textbf{Yıl} \\
\hline
Lisans & OSTİM Teknik Üniversitesi, Yazılım Mühendisliği & 2022\,--\,2026 \\
\end{tabular}

\vspace{0.8cm}

\noindent\textbf{Projeler ve Deneyimler:}

\begin{itemize}
  \item \textbf{Myora -- Yapay Zekâ Destekli Sağlık Platformu (2025--2026):}
        React 18, Supabase (PostgreSQL + Deno Edge Functions),
        OpenAI GPT-4o, RAG (\texttt{pgvector}) ve Capacitor 7 kullanılarak
        geliştirilen, 13 modüllü mobil-öncelikli sağlık yönetim platformu.
        Bitirme projesi.
\end{itemize}

\vspace{0.8cm}

\noindent\textbf{Teknik Beceriler:}

\begin{itemize}
  \item \textbf{Programlama Dilleri:} TypeScript, JavaScript, Python, SQL
  \item \textbf{Frontend:} React, Vite, Tailwind CSS, Radix UI / shadcn/ui
  \item \textbf{Backend:} Supabase (Auth, PostgreSQL, Edge Functions), Deno
  \item \textbf{Veritabanı:} PostgreSQL 15+, pgvector, Row Level Security
  \item \textbf{Yapay Zekâ:} OpenAI API (GPT-4o, text-embedding-3-small), RAG mimarisi
  \item \textbf{Mobil:} Capacitor 7 (Android / iOS)
  \item \textbf{Araçlar:} Git, Docker, VS Code, Vite, ESLint
\end{itemize}
